\documentclass[10pt,a3paper]{article}
\usepackage[margin=10mm]{geometry}
\usepackage[T1]{fontenc}
\usepackage[utf8]{inputenc}
\usepackage[french]{babel}
\usepackage{lmodern}
\usepackage{microtype}
\makeatletter\def\input@path{{../../../Style/}}\makeatother
\NeedsTeXFormat{LaTeX2e}
\ProvidesPackage{TLPosterCarteMentale}[2026/08/03 Posters et cartes mentales SI]

\RequirePackage{tikz}
\RequirePackage[most]{tcolorbox}
\RequirePackage{xcolor}
\RequirePackage{amsmath,amssymb}
\RequirePackage{graphicx}
\usetikzlibrary{positioning,arrows.meta,calc,shapes.geometric,mindmap,backgrounds,fit}

\definecolor{TLPosterBlue}{RGB}{24,86,141}
\definecolor{TLPosterLight}{RGB}{234,243,250}
\definecolor{TLPosterAccent}{RGB}{232,126,34}
\definecolor{TLPosterGreen}{RGB}{52,152,102}
\definecolor{TLPosterRed}{RGB}{192,57,43}
\definecolor{TLPosterGray}{RGB}{80,90,100}

\newcommand{\TLPosterTitre}[2]{%
  \begin{tcolorbox}[enhanced,boxrule=0pt,arc=0pt,colback=TLPosterBlue,left=6mm,right=6mm,top=4mm,bottom=4mm]
    {\color{white}\bfseries\LARGE #1}\hfill{\color{white}\large #2}
  \end{tcolorbox}%
}

\newtcolorbox{TLPosterBloc}[2][]{
  enhanced,breakable,
  colback=TLPosterLight,
  colframe=TLPosterBlue,
  colbacktitle=TLPosterBlue,
  coltitle=white,
  fonttitle=\bfseries,
  title={#2},
  boxrule=0.8pt,
  arc=2mm,
  left=3mm,right=3mm,top=2mm,bottom=2mm,
  #1
}

\newtcolorbox{TLPosterFormule}[1][]{
  enhanced,
  colback=white,
  colframe=TLPosterAccent,
  boxrule=1pt,
  arc=2mm,
  fontupper=\bfseries,
  halign=center,
  #1
}

\newcommand{\TLPosterMethode}[1]{%
  \begin{tcolorbox}[enhanced,colback=TLPosterGreen!8,colframe=TLPosterGreen,title=Méthode,fonttitle=\bfseries]
  #1
  \end{tcolorbox}%
}

\newcommand{\TLPosterAttention}[1]{%
  \begin{tcolorbox}[enhanced,colback=TLPosterRed!6,colframe=TLPosterRed,title=Point de vigilance,fonttitle=\bfseries]
  #1
  \end{tcolorbox}%
}

\tikzset{
  TLmindroot/.style={concept,concept color=TLPosterBlue,text=white,font=\bfseries\large,minimum size=34mm},
  TLmindlevel1/.style={level 1 concept/.append style={font=\bfseries,minimum size=23mm,sibling angle=45}},
  TLmindlevel2/.style={level 2 concept/.append style={font=\small,minimum size=17mm}},
  TLarrow/.style={-{Stealth[length=3mm]},thick,draw=TLPosterBlue},
  TLbox/.style={rounded corners=2mm,draw=TLPosterBlue,fill=TLPosterLight,align=center,inner sep=3mm}
}

\newenvironment{TLCarteMentale}[1]{%
  \begin{tikzpicture}[mindmap,grow cyclic,concept color=TLPosterBlue,every node/.style={concept},TLmindlevel1,TLmindlevel2]
  \node[TLmindroot]{#1}
}{;
  \end{tikzpicture}
}

\newcommand{\TLBranche}[3][]{child[concept color=#2]{node[#1]{#3}}}

\endinput

\pagestyle{empty}
\renewcommand{\familydefault}{\sfdefault}
\begin{document}
\TLPosterCourseHeader{7}{Acquisition et traitement du signal}{Choisir les outils de mesure, de filtrage et de numérisation}
\vspace{2mm}
\begin{minipage}[t]{0.60\linewidth}
\begin{TLPosterSavoirs}
\begin{itemize}
  \item choisir un capteur avec étendue, sensibilité, précision, résolution, rapidité et bande passante
  \item utiliser Kirchhoff, diviseurs, Millman, superposition et AOP pour conditionner un signal
  \item reconnaître un filtre par ses comportements basse et haute fréquence et exploiter son Bode
  \item passer du signal temporel au spectre et prévoir l'action d'un filtre sur chaque raie
  \item appliquer Shannon--Nyquist, anti-repliement, CAN/CNA, quantum et codages binaire/hexadécimal
\end{itemize}
\end{TLPosterSavoirs}
\end{minipage}\hfill
\begin{minipage}[t]{0.37\linewidth}
\TLPosterKey{Quel outil pour le circuit ?}{Série $\rightarrow$ diviseur de tension\\Parallèle $\rightarrow$ diviseur de courant\\Plusieurs sources vers un nœud $\rightarrow$ Millman\\Plusieurs sources indépendantes $\rightarrow$ superposition\\Réseau général $\rightarrow$ Kirchhoff / Kennelly}
\end{minipage}
\vfill
\begin{minipage}[t]{0.49\linewidth}
\begin{TLPosterBloc}[Capteur et conditionnement]{TLPosterBlue}
\textbf{1.} Définir mesurande et plage utile. \textbf{2.} Vérifier étendue, précision, résolution et bande passante. \textbf{3.} Adapter le signal au convertisseur : gain, offset, impédance et filtrage.\\[1mm]
AOP idéal en régime linéaire avec contre-réaction :
\[
i_+=i_-=0,\qquad u_+=u_-.
\]
Inverseur : $u_s/u_e=-R_2/R_1$ ; non-inverseur : $1+R_2/R_1$.
\end{TLPosterBloc}
\begin{TLPosterBloc}[Filtre : comment l'identifier ?]{TLPosterPurple}
Remplacer mentalement les condensateurs : \textbf{BF} $\rightarrow$ ouverts ; \textbf{HF} $\rightarrow$ courts-circuits. En déduire passe-bas, passe-haut, passe-bande ou coupe-bande. Puis établir $H(p)$, repérer $f_c$ et la pente. Un ordre $n$ donne environ $20n$ dB/décade. Le gabarit fixe bande passante, transition et atténuation minimale.
\end{TLPosterBloc}
\end{minipage}\hfill
\begin{minipage}[t]{0.49\linewidth}
\begin{TLPosterBloc}[Spectre, filtrage et échantillonnage]{TLPosterGreen}
Une raie à $nf_0$ est une harmonique ; la raie à $0$ est la valeur moyenne. À chaque fréquence :
\[
A_s(f)=|H(j2\pi f)|\,A_e(f).
\]
Avant le CAN, limiter le spectre avec un \textbf{passe-bas anti-repliement}. Shannon impose $f_e>2f_{\max}$ ; en pratique choisir une marge pour rendre la transition du filtre réalisable.
\end{TLPosterBloc}
\begin{TLPosterBloc}[CAN / CNA : dimensionner]{TLPosterTeal}
Pour un CAN sur $N$ bits et une plage $[U_{min},U_{max}]$ :
\[
q=\frac{U_{max}-U_{min}}{2^N}.
\]
Si un quantum maximal est imposé, résoudre cette relation pour trouver $N$. Le conditionneur doit exploiter la plage sans saturation. Vérifier aussi $t_{conv}<T_e=1/f_e$. Pour une valeur signée, connaître le complément à deux ; les conversions base 2--10--16 doivent être maîtrisées.
\end{TLPosterBloc}
\end{minipage}
\vfill
\begin{TLPosterMethode}
\begin{enumerate}
  \item Partir du cahier des charges : grandeur, plage, précision, fréquence maximale et dynamique attendue.
  \item Choisir le capteur puis calculer le conditionnement avec l'outil de circuit le plus direct.
  \item Si un filtre est demandé : identifier son type par BF/HF, établir $H(p)$ puis vérifier le gabarit.
  \item Si un spectre est fourni : traiter les raies une par une avec le module du filtre.
  \item Choisir $f_e$ après le filtre anti-repliement, avec $f_e>2f_{\max}$ et une marge réaliste.
  \item Dimensionner $N$, le quantum et le gain du conditionneur ; vérifier absence de saturation, repliement et perte de résolution.
\end{enumerate}
\end{TLPosterMethode}
\vfill
\TLPosterRetenir{une chaîne d'acquisition se dimensionne de la grandeur physique vers le numérique : capteur $\rightarrow$ adaptation $\rightarrow$ filtre $\rightarrow$ échantillonnage $\rightarrow$ quantification.}
\end{document}
